% =====================================================================
%  The Patient Quant — front matter
%  Half-title, title page, copyright, dedication, epigraph, contents,
%  figures, tables, management summary, preface, how to use this book.
% =====================================================================

% ---- Half-title ---------------------------------------------------------
\pagestyle{plain}
\clearpage
\thispagestyle{empty}
\vspace*{1.15in}
{\sffamily\fontsize{22}{26}\selectfont\color{txt}\bookTitle\par}
\vspace{10pt}
{\sffamily\fontsize{11}{14}\selectfont\color{txtdim}\bookSubtitle\par}
\vfill
{\sffamily\small\color{txtdim}\bookAuthor\par}
\vspace{6pt}
{\sffamily\small\color{txtdim}\bookImprint\par}
\clearpage

% ---- Title page ---------------------------------------------------------
\clearpage
\thispagestyle{empty}
\vspace*{0.75in}
{\sffamily\fontsize{26}{30}\selectfont\bfseries\color{txt}\bookTitle\par}
\vspace{12pt}
{\sffamily\fontsize{12}{16}\selectfont\color{accent}\bookSubtitle\par}
\vspace{22pt}
\noindent{\color{accent}\rule{\textwidth}{1.8pt}}\par
\vspace{18pt}
{\sffamily\large\color{txt}\bookAuthor\par}
\vspace{8pt}
{\sffamily\small\color{txtdim}\bookImprint, publisher\par}
\vspace{4pt}
{\sffamily\small\color{txtdim}\bookEdition\par}
\vspace{18pt}
{\sffamily\small\color{txtdim}
Facts current as of \bookFactsDate.\par}
\vfill
\clearpage

% ---- Copyright page -----------------------------------------------------
\clearpage
\thispagestyle{plain}
\noindent{\sffamily\bfseries\color{txt}\bookTitle}\par
\noindent{\sffamily\small\color{txtdim}\bookSubtitle}\par
\vspace{14pt}
\noindent\rule{\textwidth}{0.4pt}\par
\vspace{10pt}
\noindent Copyright \textcopyright\ 2026 Tico Internet LLC. All rights
reserved.\par
\vspace{6pt}
\noindent Tico Internet LLC, publisher.\par
\vspace{10pt}
\noindent\textbf{Edition:} 2026 First Edition.\par
\vspace{6pt}
\noindent\textbf{Subject headings:} Portfolio management --- Personal
finance; Portfolio performance --- Evaluation; Machine learning ---
Applications; Bayesian analysis; Risk management; Capital gains tax;
Mutual funds --- United States; Exchange-traded funds; Saving and
investing; Risk.\par
\vspace{10pt}
\noindent\textbf{Licence.} This work is published under the Apache License,
Version 2.0. The complete licence text is available at
\href{https://www.apache.org/licenses/LICENSE-2.0}{\code{apache.org/licenses/LICENSE-2.0}}.
The book may be copied, modified and redistributed, including for commercial
purposes, provided that this notice and the attribution required by the
licence are retained.\par
\vspace{10pt}
\noindent\textbf{DOI:} \href{https://doi.org/\bookDOI}{\texttt{\bookDOI}}\par
\vspace{10pt}
\noindent\textbf{ISBN:} 978-1-0000-0000-0 (reserved; to be assigned on
publication).\par
\vspace{10pt}
\noindent\textbf{Source materials.} The Python package, the build scripts,
the figures, the data cache and the research ledger behind this book are
published at \texttt{github.com/lordxmen2k/The-Patient-Quant}. Every figure
in the book is computed by that code from the data in that repository; no
figure is transcribed.\par
\vspace{10pt}
\noindent\textbf{What this book is.} This book builds one complete
investment system for a single investor, runs it over
\numRecordMonths{} months of history net of trading costs and of the
capital gains tax, and reports what it is worth. It carries its own
arithmetic: the estimators, the tax engine and the evaluation machinery
are written from their objective functions and are included with the
book. It states its losses plainly, and its central result is a null
result: the machine-learning layer contributes approximately nothing, and
the risk layer contributes a materially shallower worst episode. It does
not promise an edge, and it does not describe a live trading record.\par
\vspace{10pt}
\noindent\rule{\textwidth}{0.4pt}\par
\vspace{8pt}
{\sffamily\small\color{txtdim}Typeset in XeLaTeX with Outfit, DejaVu Sans Mono
and Fira Math. First edition, 2026.\par}
\clearpage

% ---- Dedication ---------------------------------------------------------
\clearpage
\thispagestyle{empty}
\vspace*{2.6in}
\begin{center}
{\sffamily\itshape\color{txtdim}
For the person who has to keep holding the statement,\\[6pt]
and who would rather know.\par}
\end{center}
\vfill
\clearpage

% ---- Epigraph -----------------------------------------------------------
\clearpage
\thispagestyle{plain}
\begin{epigraph}
The best time to plant a tree was twenty years ago. The second best time
is now, and the second best answer is to hold the index and stop reading.

\rule{1.6in}{0.3pt}

This book, on the evidence it gathers
\end{epigraph}
\vfill
\clearpage

% ---- Contents -----------------------------------------------------------
\clearpage
\thispagestyle{plain}
\markboth{Contents}{Contents}
\tableofcontents
\addcontentsline{toc}{chapter}{Contents}
\clearpage

% ---- Figures ------------------------------------------------------------
\clearpage
\thispagestyle{plain}
\markboth{Figures}{Figures}
\listoffigures
\addcontentsline{toc}{chapter}{List of Figures}
\clearpage

% ---- Tables -------------------------------------------------------------
\clearpage
\thispagestyle{plain}
\markboth{Tables}{Tables}
\listoftables
\addcontentsline{toc}{chapter}{List of Tables}
\clearpage

% ---- Management summary -------------------------------------------------
% The front matter's unnumbered chapters must not number their sections.
\setcounter{secnumdepth}{-1}
\chapter*{Management Summary}
\addcontentsline{toc}{chapter}{Management Summary}
\markboth{Management Summary}{Management Summary}
\thispagestyle{plain}

This page is the book in the form a decision-maker needs it. The remainder
of the volume derives, tests and qualifies every statement below.

\section*{The question}

Whether there is a machine-learning procedure that a single investor can
implement, that makes money and protects principal at the same time, and
that can be evaluated honestly enough that the investor knows which of the
two it is actually doing.

\section*{The answer}

A complete, specified, runnable system --- a stated strategic allocation
that the model may not touch, a bounded forecast-driven tilt inside each
sleeve, a risk governor, a capital-gains tax engine and a monthly
rebalancing loop --- evaluated over \numRecordMonths{} months, net of
spread costs and of tax, against the benchmark and against four
competitors, with the trial family declared in advance and the selection
tests applied to the book's own results.

\section*{The rule}

The strategic allocation is never set by the model. The model may tilt
inside a sleeve by a bounded amount, and it may hold cash. Everything the
model does is therefore reversible by the reader without changing the
book's argument, and the protection layer --- which needs no forecast and
therefore cannot be wrong about the future --- is what carries the result.

\section*{The five numbers that decide it}

\begin{itemize}
\item \textbf{\numFullCAGR{}} a year, against the index's
\textbf{\numSpyCAGR{}}, over \numRecordMonths{} months.
\item \textbf{\numFullDD{}}, against the index's \textbf{\numSpyDD{}} ---
a shallower worst episode by about seventeen percentage points.
\item \textbf{\numFullSharpe{}}, against the index's \textbf{\numSpySharpe{}}
--- indistinguishable on risk-adjusted return.
\item \textbf{\numPBO{}} --- the probability that the best configuration in
the first half of the record is among the worst in the second half is
above one half. The selection procedure is worse than a coin flip.
\item \textbf{\numTiltMean{}} a year --- what the machine-learning layer
contributes. Approximately nothing, in both senses of the phrase.
\end{itemize}

\section*{What it does not do}

It does not beat the index. It does not claim that any forecast was
found. It does not argue that the equal-weighted portfolio of the same
sixty-one funds is worse --- on this record that portfolio earned
\textbf{\numEqualCAGR{}} with a drawdown of \textbf{\numEqualDD{}}, more
money and more pain, with no machinery at all. And it does not describe a
live trading record; it describes a simulation under stated assumptions.

\begin{keytakeaways}
\item The protection comes from the risk layer, which needs no forecast,
and its contribution is positive in every starting month tested.
\item The prediction layer contributes approximately nothing, and the
selection tests cannot distinguish the system's return from the
benchmark's.
\item The evaluation is the transferable result: a declared trial family
of \numFamilySize{} members, four selection tests applied to the book's
own results, and a null reported as a null.
\item The largest number in the subject is behavioural: the psychic cost of
evaluating a long-horizon portfolio annually is 5.1 per cent a year.
\end{keytakeaways}
\clearpage

% ---- Preface ------------------------------------------------------------
\chapter*{Preface}
\addcontentsline{toc}{chapter}{Preface}
\markboth{Preface}{Preface}
\thispagestyle{plain}

The owner of this project asked for a book about making money and
protecting principal with quant prediction algorithms and machine learning,
and the request was framed in the language of the trading floor: getting
every advantage and avoiding every pitfall.

Written to that specification, the book has an obvious and wrong ending.
A system that predicts returns, trades them, and reports a Sharpe ratio
would satisfy the request and would also be worth nothing, for reasons
that are now well established and that the first four chapters set out.
The replication evidence is partial and the published factors have decayed.
Most professional managers, with orders of magnitude more data and more
people, lose to a passive index on every long window measured. A third of
the published equity-premium predictors no longer work in their own
original sample.

So the book was built the other way round. The forecast was made to be
\emph{priced}, not to be believed: the model's output is a posterior
distribution over the forward return, and the portfolio layer's job is to
size positions against the width of that distribution. The measurement
then came out where the measurement was always going to come out. The
model's contribution is \numTiltMean{} a year, with a median of
\numTiltMedian{}, and it is positive in \numTiltPosFrac{} of the
windows tested. The risk governor's contribution is positive in every
one.

\section*{What the book does}

It builds the whole thing. The data layer fetches and checksums sixty-one
funds' price and dividend histories and rebuilds the monthly total returns
from their components, masking the period before each fund listed so that
a new listing does not acquire an artificial history of zero volatility.
The model layer is a conjugate normal-inverse-gamma panel regression,
refitted at every formation date on rows whose outcome has already been
realised, combined across \numModelConfigs{} configurations by an online
Bayesian model average whose weights depend only on outcomes already
known. The portfolio layer applies a stated allocation, a
certainty-equivalent Kelly tilt bounded inside each sleeve, a damped
volatility scaler and a two-level drawdown brake. The tax engine keeps
lots, applies the wash-sale rule in both directions across two accounts,
and settles annually. The evaluation layer computes the compound growth
rate, the Lo-corrected Sharpe ratio, the maximum drawdown, the time to
recovery, the worst rolling twelve-month outcome, the Newey--West alpha,
the Deflated Sharpe Ratio, the probability of backtest overfitting, the
White reality check and the Hansen superior predictive ability test, and
applies them all to the book's own family.

None of it is a black box. Every estimator is written from its objective
function, including the normal and Student distributions, and the package
runs with nothing but \code{numpy}, \code{pandas} and \code{matplotlib}.

\section*{What the book concludes}

That the prediction is the easy part and it does not work; that the
protection is the hard part and it does; and that the architecture which
lets the two coexist --- a stated allocation the machine may not touch, a
bounded tilt, a posterior that prices its own uncertainty, a governor
that needs no forecast, a tax engine, and a trial family declared in
advance --- is the result worth carrying away.

\section*{How to read it if you are in a hurry}

Read the Management Summary, then Part VII. If you want to know whether
the numbers are real, read Appendix A and the verification gates in
Appendix D, and then run the code.

\begin{keytakeaways}
\item The book was written to a specification whose obvious ending would
have been worth nothing, so it was built the other way round: the
forecast is priced rather than believed.
\item Every estimator is written from its objective function; the package
runs on \code{numpy}, \code{pandas} and \code{matplotlib} alone.
\item The conclusion is that the risk layer carries the result and the
prediction layer does not, and that the architecture is the transferable
part.
\end{keytakeaways}
\clearpage

% ---- How to use this book ------------------------------------------------
\chapter*{How to Use This Book}
\addcontentsline{toc}{chapter}{How to Use This Book}
\markboth{How to Use This Book}{How to Use This Book}
\thispagestyle{plain}

\section*{If you intend to run the system}

Start with Appendix A, which documents the data and its provenance, and
Appendix C, which lists every fund, its expense ratio, its inception date
and the live window over which the book could actually hold it. Then read
Chapter 30, which describes the monthly loop in enough detail to
reproduce it, and Chapter 33, which reports what the system is worth
after every cost.

The package is at \texttt{github.com/lordxmen2k/The-Patient-Quant}. The
data fetch is a single command; the analysis is a single command; the
figures are a single command. The build script runs the whole chain.

\section*{If you intend to argue with it}

The chapters most likely to be argued with are Chapter 23, which fixes the
strategic allocation outside the model, Chapter 25, which reports the
evidence that optimised portfolios do not beat equal weighting, and
Chapter 32, which reports that the equal-weighted portfolio beat this
system on compound return over the record.

The book's own numbers are in \texttt{build/numbers.json} and in the
generated macro file, and every figure is computed from those files. If a
number in the text disagrees with the code, the code is right and the text
is wrong, and the verification gates are set up to catch exactly that.

\section*{If you are only reading for the argument}

Part I states the question and the evidence that bounds it. Part VII gives
the verdict and the recommendations. Part IV contains the mathematics and
is the part that will matter if you want to build one of these yourself.

\begin{keytakeaways}
\item Run it: Appendix A and Appendix C document the data, Chapter 30 the
loop, and the repository runs the whole chain with one command.
\item Argue with it: Chapters 23, 25 and 32 are the load-bearing claims,
and every number behind them is in the repository.
\item Read it for the argument: Parts I and VII, with Part IV if you intend
to build one of these yourself.
\end{keytakeaways}
\clearpage
